\section{Bounded canonical potentials and nefness}
\label{sec:nefness}

The first application of the maximizing inequality proves nefness of
$K_X$. We proceed in three steps: bound the Poisson Ricci integral on
every finite-area disc,
convert the bound into locally bounded plurisubharmonic metric weights,
and smooth those weights with arbitrarily small loss of positivity.
We first apply the preceding disc inequalities with the fixed K\"ahler
form $k$.  As before, $R_k=-\ddc\log\det k$, and $D$ denotes the
uniform upper bound for $d(f)=k(f'(0),f'(0))$ from
Lemma~\ref{ex:brody}.

\begin{proposition}[Uniform Ricci-disc bound]
\label{nef:disc-bound}
There is a constant $C\geq 0$, depending only on the fixed metric $k$,
such that
\begin{equation}
 P_{R_k}(f)\leq C
 \label{eq:nef-uniform-disc-bound}
\end{equation}
for every finite-area holomorphic disc $f:\D\to X$.
\end{proposition}

\begin{proof}
For $0<\delta\leq 1$ and a prescribed center $x$, let $f$ maximize
\[
J_\delta(f)=d(f)+P_{R_k}(f)-\delta A_k(f).
\]
Such a maximizer exists by Proposition~\ref{ex:maximizer}.
Use the boundary-normalized holomorphic frame $B$ associated to this
maximizer and put $N(z)=B(z)^*k(f(z))B(z)$.  Set
\[
 M=\sup_X|\tr_k R_k|.
\]
The dilation inequality, Equation~\eqref{eq:tail}, gives
$A_{R_k}(f)\geq-d(f)\geq-D$.
Consequently, Equation~\eqref{eq:maximum-inequality}, with $h=k$ and
$q=R_k$, implies
\[
 \tr N(0)
 \leq \delta\bigl(n-A_{R_k}(f)\bigr)+M
 \leq n+D+M.
\]
The determinant identity, Equation~\eqref{eq:determinant-identity},
and the arithmetic-geometric mean inequality yield
\[
 P_{R_k}(f)=\log\det N(0)
 \leq n\log\!\left(\frac{n+D+M}{n}\right)=:C_1.
\]
Since $A_k(f)\geq0$ and $d(f)\leq D$, the maximal value of
$J_\delta$ is at most $D+C_1$, independently of both $x$ and $\delta$.
Now fix any finite-area disc $g$, use its center in this maximization,
and obtain
\[
 d(g)+P_{R_k}(g)-\delta A_k(g)\leq D+C_1.
\]
Letting $\delta\downarrow0$ for this fixed disc and using $d(g)\geq0$
proves the assertion with $C=D+C_1$.  This limiting argument requires
no area bound uniform in $\delta$ for the maximizing discs.
\end{proof}

\subsection{The Poisson envelope on the determinant bundle}

Poletsky introduced disc envelopes for plurisubharmonic functions,
and Rosay established the manifold case \cite{Poletsky1991,Rosay2003}.
The formulation below is due to Drinovec Drnov\v sek and Forstneri\v c,
whose theorem also covers locally irreducible complex spaces
\cite{DrinovecForstneric2012}. The use of a line bundle lets the
Ricci-disc bound control these envelopes without choosing one global
potential on $X$.

We use the Poletsky--Rosay Theorem in the following form.  If $Z$ is
a connected complex manifold and $b:Z\to\R\cup\{-\infty\}$ is upper
semicontinuous, then
\[
 \mathcal P b(y)
 =\inf_{G(0)=y}\frac{1}{2\pi}\int_0^{2\pi}b(G(e^{i\theta}))\,\dd\theta
\]
is plurisubharmonic or identically $-\infty$.  Here the infimum is
taken over maps $G:\overline{\D}\to Z$ that are continuous on
$\overline{\D}$ and holomorphic in $\D$.  There is no requirement that
$Z$ be compact or that $b$ be bounded below.  This is, in particular,
the manifold case of the more general Theorem~1.1 of
\cite{DrinovecForstneric2012}.  For continuous real-valued $b$, the
same envelope results if the test maps are required to be holomorphic
on neighborhoods of $\overline{\D}$: replace each test map $G$ by
$z\mapsto G(rz)$ and let $r\uparrow1$.

\begin{proposition}[Bounded positive canonical potentials]
\label{nef:bounded-potentials}
There is a bounded real-valued function $u$ on $X$ such that
\[
 T=-R_k+\ddc u\geq0
\]
is a closed positive $(1,1)$-current.  In every local holomorphic
frame $s_i$ of $\det T^{1,0}X$, put $a_i=\log|s_i|_k^2$.
Then $a_i+u$ is plurisubharmonic and locally bounded.  These are
local weights of a singular Hermitian metric with nonnegative
curvature on $K_X$.
\end{proposition}

\begin{proof}
Write $L=\det T^{1,0}X$ and let $\pi:Y=L\setminus 0_X\to X$ be
the complement of its zero section.  Define
\[
 \ell(s)=\log|s|_k^2,\qquad s\in Y.
\]
This is a smooth finite-valued function on the connected complex
manifold $Y$.  In a local trivialization $s=\lambda s_i(x)$ it has
the expression
\[
 \ell(s)=\log|\lambda|^2+a_i(x),\qquad
 \ddc a_i=-R_k.
\]
Since $\lambda\ne0$, the first summand is pluriharmonic.  Therefore
\begin{equation}
 \ddc\ell=-\pi^*R_k.
 \label{eq:nef-total-space-curvature}
\end{equation}

Let $v=\mathcal P\ell$, using test discs holomorphic past the closed
unit disc.  For such a disc $G$ with $G(0)=s$, its projection
$f=\pi\circ G$ has finite area.  Green's formula and
Equation~\eqref{eq:nef-total-space-curvature} give
\[
 \frac{1}{2\pi}\int_0^{2\pi}\ell(G(e^{i\theta}))\,\dd\theta
 =\ell(s)-P_{R_k}(f).
\]
Proposition~\ref{nef:disc-bound} bounds every such mean below by
$\ell(s)-C$, while the constant disc has mean $\ell(s)$.  Thus
\begin{equation}
 \ell-C\leq v\leq\ell.
 \label{eq:nef-envelope-bounds}
\end{equation}
The Poletsky--Rosay Theorem applies even though $\ell$ is globally
unbounded, and Equation~\eqref{eq:nef-envelope-bounds} excludes its
identically $-\infty$ alternative.  Hence $v$ is plurisubharmonic.

For every $\lambda\in\C^*$, multiplication by $\lambda$ is a
biholomorphism of $Y$ and gives a bijection between the test discs
centered at $s$ and at $\lambda s$.  Since
$\ell(\lambda s)=\ell(s)+\log|\lambda|^2$, it follows that
\[
 v(\lambda s)=v(s)+\log|\lambda|^2.
\]
Thus $v-\ell$ is constant on every fiber.  It has the form
$u\circ\pi$, where $-C\leq u\leq0$.  Pulling back $v$ by the local
holomorphic map $s_i$ proves that $a_i+u$ is plurisubharmonic.
It is locally bounded because $a_i$ is smooth and $u$ is bounded.
The local currents $\ddc(a_i+u)$ agree and equal $-R_k+\ddc u$.

To check the line-bundle sign explicitly, if $s_i=g_{ij}s_j$ then
\[
 a_i-a_j=\log|g_{ij}|^2.
\]
Let $s_i^*$ be the dual frame of $K_X=L^*$.  The prescription
\[
 |s_i^*|^2=\exp\bigl(-(a_i+u)\bigr)
\]
is compatible with $s_i^*=g_{ij}^{-1}s_j^*$.  Its curvature is
$\ddc(a_i+u)=T\geq0$.  In particular $T$ represents
$2\pi c_1(K_X)$, rather than $2\pi c_1(L)$.
\end{proof}

\subsection{From bounded potentials to nefness}

A holomorphic line bundle is nef if, for every $\eps>0$, it admits a
smooth Hermitian metric whose real curvature is at least $-\eps k$.
This definition is independent of the background K\"ahler metric on
a compact manifold. The following approximation result converts the
bounded potentials just constructed into such metrics.

\begin{lemma}[Smoothing a bounded positive potential]
\label{nef:bounded-smoothing}
Let $(M,k)$ be a compact K\"ahler manifold, let $\alpha$ be a smooth
closed real $(1,1)$-form on $M$, and let $u$ be a bounded real-valued
function such that $\alpha+\ddc u\geq0$ in the sense of currents.
For every $\eps>0$ there is a smooth real-valued function $u_\eps$
on $M$ such that
\[
 \alpha+\ddc u_\eps\geq-\eps k.
\]
\end{lemma}

This is the bounded-potential case of the regularization theorem for
positive currents: locally bounded plurisubharmonic potentials have
zero Lelong numbers, so \cite[Corollary~6.4]{Demailly1992} gives
nefness. Appendix~\ref{app:smoothing} supplies a direct proof. Its
key estimate compares regularizations on overlapping charts without
assuming that the original bounded potentials are continuous.

\begin{proposition}[Nefness of the canonical class]
\label{nef:nef}
For every $\eps>0$ there is a smooth real-valued function $u_\eps$
on $X$ such that
\[
 -R_k+\ddc u_\eps\geq-\eps k.
\]
Consequently $2\pi c_1(K_X)$ is nef.
\end{proposition}

\begin{proof}
Apply Lemma~\ref{nef:bounded-smoothing} with $M=X$, $\alpha=-R_k$,
and the bounded function $u$ of Proposition~\ref{nef:bounded-potentials}.
The resulting functions $u_\eps$ give smooth canonical metrics with
local squared norms $\exp(-(a_i+u_\eps))$ in the dual frames $s_i^*$.
Their curvatures are $-R_k+\ddc u_\eps\geq-\eps k$.
\end{proof}
